\section{Aplikasi dan Studi Kasus}

Bagian ini menyajikan dua studi kasus yang mengaitkan ekuivalensi logika, hukum De Morgan, serta tautologi dan kontradiksi dengan pemrograman dan verifikasi aturan bisnis \cite{levin_discrete,python_docs_boolean,rosen2019}.

\subsection*{Studi Kasus 1: Penyederhanaan Kondisi Program dengan Hukum De Morgan}

Dalam pemrograman, kondisi percabangan sering memuat negasi dari konjungsi atau disjungsi. Hukum De Morgan memungkinkan kita menulis ulang kondisi tersebut menjadi bentuk yang kadang lebih mudah dibaca atau lebih efisien. Misalnya: ``akses ditolak jika dan hanya jika pengguna \emph{tidak} (terdaftar \emph{dan} kata sandi benar)''. Bentuk asli: $\neg(p \wedge q)$ dengan $p$ = ``terdaftar'', $q$ = ``kata sandi benar''. Menurut De Morgan: $\neg(p \wedge q) \equiv \neg p \vee \neg q$, yaitu ``tidak terdaftar \emph{atau} kata sandi salah''. Kedua bentuk ekuivalen; kita bisa memilih yang lebih jelas untuk konteks kode \cite{brilliant_de_morgan,rosen2019}.

\begin{contoh}
\textbf{Sebelum dan sesudah penerapan De Morgan dalam Python}

Kondisi ``jangan izinkan jika (sudah login \textbf{dan} bukan admin)'' ekuivalen dengan ``izinkan jika (belum login \textbf{atau} admin)''.

\begin{lstlisting}[style=pythonstyle, caption={Penyederhanaan kondisi dengan De Morgan: tidak (A dan B) = tidak A atau tidak B}]
def boleh_akses_v1(sudah_login, is_admin):
    """Versi asli: akses ditolak jika (login dan bukan admin)."""
    akses_ditolak = sudah_login and (not is_admin)
    return not akses_ditolak

def boleh_akses_v2(sudah_login, is_admin):
    """Versi ekuivalen (De Morgan): izinkan jika tidak login ATAU admin."""
    return (not sudah_login) or is_admin

# Kedua versi memberi hasil sama untuk setiap (sudah_login, is_admin)
for login, admin in [(True, True), (True, False), (False, True), (False, False)]:
    assert boleh_akses_v1(login, admin) == boleh_akses_v2(login, admin)
print("De Morgan: kedua bentuk ekuivalen.")
\end{lstlisting}
\end{contoh}

\subsection*{Studi Kasus 2: Verifikasi Aturan Bisnis dengan Tautologi dan Implikasi}

Aturan bisnis sering dinyatakan sebagai implikasi: ``jika mahasiswa lulus maka dia sudah membayar SPP.'' Misalkan $p$ = ``Mahasiswa lulus'', $q$ = ``Sudah bayar SPP''. Spesifikasi: tidak pernah terjadi (lulus dan belum bayar), yaitu $\neg(p \wedge \neg q)$, yang ekuivalen dengan $p \rightarrow q$. Untuk memverifikasi bahwa aturan konsisten, kita dapat memeriksa bahwa implikasi $p \rightarrow q$ bersifat kontingen (bukan kontradiksi); jika ada baris yang membuat $p \rightarrow q$ salah, aturan tersebut dilanggar. Dengan Python kita bisa mengevaluasi $p \rightarrow q$ untuk semua interpretasi dan memastikan tidak ada kasus ``lulus tapi belum bayar'' (yaitu $p$ benar dan $q$ salah) dalam data yang valid \cite{levin_discrete,rosen2019}.

\begin{contoh}
\textbf{Memverifikasi aturan ``lulus $\rightarrow$ sudah bayar'' dan cek kontradiksi}

\begin{lstlisting}[style=pythonstyle, caption={Verifikasi aturan bisnis: implikasi dan deteksi pelanggaran}]
from itertools import product

def impl(p, q):
    """p -> q"""
    return (not p) or q

# Aturan: jika lulus maka sudah bayar (p -> q)
# Pelanggaran: p True dan q False
print("p (lulus)  q (bayar)  p->q   Pelanggaran?")
print("-" * 50)
for p, q in product([True, False], repeat=2):
    valid = impl(p, q)
    violation = p and (not q)
    pl = "B" if p else "S"
    ql = "B" if q else "S"
    vl = "B" if valid else "S"
    vstr = "Ya" if violation else "Tidak"
    print(f"{pl:10} {ql:10} {vl:5} {vstr}")

# Cek: (p -> q) <-> (~p or q) tautologi?
def is_tautology_impl_equiv():
    for p, q in product([True, False], repeat=2):
        if impl(p, q) != ((not p) or q):
            return False
    return True
print("\n(p -> q) ekuivalen (~p or q) untuk semua (p,q):", is_tautology_impl_equiv())
\end{lstlisting}
\end{contoh}

Kedua studi kasus di atas menunjukkan penerapan ekuivalensi logika dan hukum De Morgan dalam penyederhanaan kondisi program, serta penggunaan tautologi/implikasi untuk memverifikasi aturan bisnis \cite{python_docs_boolean,rosen2019}.
